\section{Introducción}
\label{sec:introduction}

El transporte de una carga que excede a un solo robot obliga a decidir quién participa y si el equipo puede ejecutar la tarea. Ambas preguntas están acopladas: una asignación válida puede fallar por geometría, torque, saturación, colisión o pérdida de un integrante.

\subsection{Contexto}

La asignación multi-robot está bien delimitada cuando cada tarea corresponde a un robot \parencite{gerkeyMataric2004MRTA}. Las cargas voluminosas, frágiles o pesadas rompen esa estructura. Requieren varios robots al mismo tiempo, combinan capacidades y crean dependencias entre ganadores \parencite{korsahStentzDias2013iTax}. Las coaliciones temporales permiten adaptar la flota a cada carga sin dimensionar todos los robots para el peor caso \parencite{shehoryKraus1998Coalition,vigAdams2006Coalition}.

La relevancia técnica no reside solo en repartir tareas. Una decisión de
reclutamiento modifica la geometría del conjunto móvil, las fuerzas disponibles,
el espacio ocupado, el tráfico y la capacidad de recuperarse. Evaluar cada capa
por separado puede aceptar una coalición que la capa siguiente no puede
ejecutar. Por ello, el problema científico consiste en definir interfaces
verificables entre decisión, factibilidad física y movimiento, además de medir
qué se pierde al sustituir información global por comunicación vecinal.

Contar integrantes no basta. El equipo debe cerrar su composición, ocupar contactos compatibles y producir la fuerza y el torque requeridos. ASyMTRe e IQ-ASyMTRe relacionan capacidades funcionales con composiciones ejecutables \parencite{parkerTang2006ASyMTRe,zhangParker2013IQASyMTRe}; la teoría de agarre muestra que la geometría y el contacto determinan el cierre y la distribución de esfuerzos \parencite{bicchi1995closure,nakamura1989dynamics}. La asignación y la ejecución física necesitan certificados distintos.

\subsection{Problema de investigación}

Se considera una flota heterogénea y un conjunto de cargas con requisitos de capacidad, contacto y pose. La arquitectura objetivo privilegia estado propio, percepción local y mensajes vecinales, pero la implementación evaluada es híbrida: combina actualizaciones locales con cierres, agregados, guardias por carga y un registro de atributos que conservan dependencias globales declaradas. La misión comprende reclutamiento, aproximación, contacto, transporte y recuperación tras fallos. Los optimizadores globales se reservan como oráculos experimentales o componentes explícitos de certificación.

El problema contiene cinco objetos que no deben confundirse. La estrategia determina preferencias y pertenencia. La estimación aproxima ocupación y recursos no observados. La planta describe robots y cargas. El contacto convierte esfuerzos individuales en \emph{wrench}. La seguridad restringe el movimiento y la recuperación. Cada capa responde solo de su propia pregunta: el equilibrio pertenece a la estrategia, la formación geométrica no acredita por sí misma la fuerza requerida y la ausencia de colisiones en simulación no constituye una prueba general de seguridad.

Las subastas resuelven conflictos de adjudicación \parencite{gerkeyMataric2002MURDOCH,choiBrunetHow2009CBBA}; los juegos de coalición modelan el valor del grupo \parencite{vigAdams2006Coalition,chenSun2012Coalition}; el control cooperativo mueve equipos ya constituidos \parencite{alonsomora2017transport}. En el corpus verificado no se identificó un método único que cubra las transiciones entre preferencia, decisión entera, factibilidad mecánica, movimiento, reparación y red imperfecta. Esta observación delimita la brecha revisada; no afirma inexistencia universal.

Los juegos potenciales modelan los incentivos; el cierre entero asigna robots identificables; los certificados físicos descartan coaliciones que no pueden ejecutar la tarea. La implementación es digital y muestreada. \emph{Asíncrono} significa que no hay rondas globales de decisión, aunque cada robot actualiza su estado con un reloj digital.

\subsection{Contribución y alcance}

La contribución consiste en una arquitectura por capas que hace explícita la cadena entre preferencia, cierre, guardia y ejecución. SP0--SP2 caracterizan asignación, cuotas y servicio heterogéneo. SP3 sustituye el índice nominal por factibilidad planar de \emph{wrench}, y SP4 relaciona el residual mecánico con el transporte de pose. SP5 filtra la huella compuesta ante obstáculos; SP6 recompone un certificado aditivo tras un fallo; SP7 coordina tráfico lógico sobre un grafo discreto y SP8 mide escala y comunicación imperfecta sobre ese mismo modelo de rutas. La simulación Cargo integra las interfaces de SP2--SP6, pero no transfiere automáticamente las garantías de un SP a otro.

La modalidad física desarrollada es Cargo nominal: después del acoplamiento, una carga y los robots se representan mediante una planta planar reducida con contactos fijos y \emph{wrench} acotado. El modelo no valida soporte vertical ni contacto rueda--suelo durante el transporte. El empuje y el \emph{caging} quedan fuera porque exigen restricciones de contacto diferentes. Tampoco se demuestra convergencia del sistema completo, estabilidad tridimensional, robustez ante cualquier perturbación ni escalabilidad ilimitada. La evidencia es formal o experimental según el SP y no se suma como una garantía extremo a extremo.

La separación localiza el origen de los fallos. Una coalición cerrada puede ser mecánicamente inviable. Un control que evita colisiones puede bloquearse. Una red local puede estabilizar decisiones construidas con información incompleta. Cada fenómeno conserva así su propia medida: factibilidad, error de pose, residual de \emph{wrench}, colisión, recuperación, comunicación o coste computacional.

\subsection{Preguntas de investigación}

La pregunta principal es si una coordinación basada en juegos y control local puede formar coaliciones factibles y transportar cargas heterogéneas con garantías delimitadas y calidad medible frente a referencias centralizadas y distribuidas.

\begin{enumerate}[label=\textbf{RQ\arabic*.},leftmargin=3.2em,topsep=0.35em,itemsep=0.25em,parsep=0pt]
  \item \parbox[t]{\linewidth}{¿Bajo qué condiciones de demanda, heterogeneidad, recursos y parámetros emergen coaliciones sin déficit ni sobrerreclutamiento persistente?}
  \item \parbox[t]{\linewidth}{¿Cómo afectan el acoplamiento entre decisión y movimiento y la comunicación local al equilibrio, al tiempo y a la calidad de la misión?}
  \item \parbox[t]{\linewidth}{¿Qué restricciones de fuerza, torque y contacto separan capacidad nominal de ejecución física?}
  \item \parbox[t]{\linewidth}{¿Cómo se degrada la operación ante congestión, batería limitada, retardo, pérdida de paquetes y fallos parciales?}
  \item \parbox[t]{\linewidth}{¿Qué coste computacional, comunicativo y operacional introduce la descentralización?}
\end{enumerate}

\subsection{Estructura de la memoria}

El Capítulo~2 convierte estas preguntas en un objetivo general y cinco objetivos
específicos; el Capítulo~3 fija las hipótesis y sus criterios de contradicción.
El Capítulo~4 describe el diseño, las unidades de análisis, los modelos, los
comparadores y la estadística. El Capítulo~5 sitúa la propuesta frente al corpus
verificado. El Capítulo~6 presenta SP0--SP8, la campaña integrada Cargo, el
piloto Industrial~2 y la discusión transversal. El Capítulo~7 responde a cada
objetivo, pregunta e hipótesis, y separa resultados sustentados, resultados
negativos, limitaciones y trabajo futuro.
